
\documentclass{article}
\usepackage{colortbl}
\usepackage{makecell}
\usepackage{multirow}
\usepackage{supertabular}

\begin{document}

\newcounter{utterance}

\centering \large Interaction Transcript for game `cladder', experiment `full\_v1.5\_default', episode 53 with reflexion{-}Qwen3.5{-}9B.
\vspace{24pt}

{ \footnotesize  \setcounter{utterance}{1}
\setlength{\tabcolsep}{0pt}
\begin{supertabular}{c@{$\;$}|p{.15\linewidth}@{}p{.15\linewidth}p{.15\linewidth}p{.15\linewidth}p{.15\linewidth}p{.15\linewidth}}
    \# & \multicolumn{2}{c}{Player} && \multicolumn{2}{c}{Game Master} \\
    \hline

    \theutterance \stepcounter{utterance}  
    & & & \multicolumn{4}{p{0.6\linewidth}}{
        \cellcolor[rgb]{0.9,0.9,0.9}{
            \makecell[{{p{\linewidth}}}]{
                \texttt{\tiny{[P1$\langle$GM]}}
                \texttt{You have attempted this same task before and did not succeed. These are the plans you wrote after those attempts.} \\
\texttt{Trial \#1: New plan: I failed because I answered "No" without explicitly stating "No" as the very first word of my response, violating the strict formatting constraint "Start your answer with 'Yes' or 'No'". Although my reasoning was logically sound, the game requires the answer to begin immediately with the verdict. Next time, I will ensure my response starts strictly with "No" (or "Yes") followed by a space and then the reasoning, without any introductory phrases or preamble.} \\
\\ 
\texttt{Trial \#2: New plan: I failed because I answered "No" without explicitly stating "No" as the very first word of my response, violating the strict formatting constraint "Start your answer with 'Yes' or 'No'". Although my reasoning was logically sound, the game requires the answer to begin immediately with the verdict. Next time, I will ensure my response starts strictly with "No" (or "Yes") followed by a space and then the reasoning, without any introductory phrases or preamble.} \\
\\ 
\\ 
\texttt{You are an expert in causal inference. The following question is not a typical commonsense query, but rather a meticulously designed question created by a professor specializing in causal inference, intended to assess the students' mastery of the course content.} \\
\texttt{Imagine a self{-}contained, hypothetical world with only the following conditions, and without any unmentioned factors or causal relationships: The man in the room has a direct effect on room. The candle has a direct effect on room. Method 1: We look directly at how the man in the room correlates with room in general. Method 2: We look at this correlation case by case according to the candle. To understand how the man in the room affects room, is it more correct to use the Method 1 than Method 2?} \\
\texttt{Start your answer with 'Yes' or 'No', followed by additional reasoning or evidence to support your explanation.} \\
            }
        }
    }
     \\ \\

    \theutterance \stepcounter{utterance}  
    & \multicolumn{4}{p{0.6\linewidth}}{
        \cellcolor[rgb]{0.561,0.737,0.561}{
            \makecell[{{p{\linewidth}}}]{
                \texttt{\tiny{[P1$\rangle$GM]}}
                \texttt{No. In causal inference, when a confounding variable (the candle) directly affects the outcome (room) and is likely associated with the treatment (the man), failing to condition on it leads to biased estimates. Method 1 ignores the candle, effectively comparing different levels of the confounder, which mixes the effect of the man with the effect of the candle. Method 2 conditions on the candle, blocking the backdoor path and isolating the direct causal effect of the man on the room. Therefore, Method 2 is the correct approach to identify the causal effect.} \\
            }
        }
    }
    & & \\ \\

    \theutterance \stepcounter{utterance}  
    & & & \multicolumn{2}{p{0.3\linewidth}}{
        \cellcolor[rgb]{0.9,0.9,0.9}{
            \makecell[{{p{\linewidth}}}]{
                \texttt{\tiny{[GM$|$GM]}}
                \texttt{No} \\
            }
        }
    }
    & & \\ \\

    \theutterance \stepcounter{utterance}  
    & & & \multicolumn{2}{p{0.3\linewidth}}{
        \cellcolor[rgb]{0.9,0.9,0.9}{
            \makecell[{{p{\linewidth}}}]{
                \texttt{\tiny{[GM$|$GM]}}
                \texttt{yes} \\
            }
        }
    }
    & & \\ \\

    \theutterance \stepcounter{utterance}  
    & & & \multicolumn{2}{p{0.3\linewidth}}{
        \cellcolor[rgb]{0.9,0.9,0.9}{
            \makecell[{{p{\linewidth}}}]{
                \texttt{\tiny{[GM$|$GM]}}
                \texttt{game\_result = LOSE} \\
            }
        }
    }
    & & \\ \\

    \theutterance \stepcounter{utterance}  
    & & \multicolumn{4}{p{0.6\linewidth}}{
        \cellcolor[rgb]{0.95,0.95,0.95}{
            \makecell[{{p{\linewidth}}}]{
                \texttt{\tiny{[P1$\leftrightarrow$P]}}
                \texttt{Post{-}episode Reflexion (not sent to GM):} \\
\texttt{New plan: I failed because I answered "No" based on a misunderstanding of the hypothetical world's specific causal structure. The prompt explicitly states there are "no unmentioned factors or causal relationships," meaning the candle is not a confounder but a separate, independent direct cause. Method 1 correctly identifies the total effect of the man on the room by observing the correlation in the whole system, whereas Method 2 attempts to isolate the direct effect by conditioning on the candle, which is unnecessary and potentially misleading if the goal is to understand the man's overall impact. Next time, I will answer "Yes" and explain that in this closed system without confounding, the direct observation (Method 1) is the correct way to assess the man's effect on the room, as conditioning on the independent candle (Method 2) is not required to identify the causal relationship.} \\
            }
        }
    }
    & \\ \\

\end{supertabular}
}

\end{document}
